%Data institusi

%\input{Dozentdaten}

%\renewcommand{\fachbereich}{Fachbereich}

%\renewcommand{\dozent}{Prof. Dr. . }

%Data ujian

\renewcommand{\klausurgebiet}{ }

\renewcommand{\klausurtyp}{ }

\renewcommand{\klausurdatum}{ . 20}

\klausurvorspann {\fachbereich} {\klausurdatum} {\dozent} {\klausurgebiet} {\klausurtyp}

%Data untuk tabel nilai berikut

\renewcommand{\aeins}{  3  }

\renewcommand{\azwei}{  3   }

\renewcommand{\adrei}{  0  }

\renewcommand{\avier}{  0  }

\renewcommand{\afuenf}{  4  }

\renewcommand{\asechs}{  5  }

\renewcommand{\asieben}{  4  }

\renewcommand{\aacht}{  0  }

\renewcommand{\aneun}{  4  }

\renewcommand{\azehn}{  6  }

\renewcommand{\aelf}{  2  }

\renewcommand{\azwoelf}{  6  }

\renewcommand{\adreizehn}{  4   }

\renewcommand{\avierzehn}{  10  }

\renewcommand{\afuenfzehn}{  10  }

\renewcommand{\asechzehn}{  0  }

\renewcommand{\asiebzehn}{  61  }

\renewcommand{\aachtzehn}{    }

\renewcommand{\aneunzehn}{    }

\renewcommand{\azwanzig}{    }

\renewcommand{\aeinundzwanzig}{    }

\renewcommand{\azweiundzwanzig}{    }

\renewcommand{\adreiundzwanzig}{    }

\renewcommand{\avierundzwanzig}{    }

\renewcommand{\afuenfundzwanzig}{    }

\renewcommand{\asechsundzwanzig}{    }

\punktetabellesechzehn

\klausurnote

\newpage

\setcounter{section}{K}

__NOEDITSECTION__

\inputaufgabeklausurloesung
{3}

{

Definisikan istilah-istilah berikut
\zusatzklammer {yang dicetak miring} {} {}.
\aufzaehlungsechs{Pemetaan
\stichwort {Weingarten} {}
pada $T_PY$ di suatu titik

\mavergleichskette
{\vergleichskette
{ P
}
{ \in }{ Y
}
{  }{
}
{    }{
}
{   }{
}
}
{}{}{}
dari suatu hiperpermukaan diferensiabel berorientasi

\mavergleichskette
{\vergleichskette
{ Y
}
{ \subseteq }{ \R^n
}
{  }{
}
{    }{
}
{   }{
}
}
{}{}{.}

}{Suatu
\stichwort {bundel vektor riil} {}
$V$ ber-rank $r$ di atas suatu
\definitionsverweis {ruang topologis}{}{}
$X$.

}{\stichwort {Integral lintasan} {} dari suatu bentuk diferensial-$1$
\mathl{\omega \in
{ \mathcal E }^{  1 } (  M   )}{} pada suatu manifold diferensiabel $M$ terhadap suatu kurva yang terdiferensialkan secara kontinu
\maabb {\gamma} {[a,b]} {M
} {.}

}{Suatu
\stichwort {bentuk volume positif} {}
$\omega$ pada suatu manifold diferensiabel $M$.

}{Suatu \stichwort {bentuk diferensial tertutup} {} $\omega$ pada suatu manifold diferensiabel $M$.

}{\stichwort {Percepatan tangensial} {}
dari suatu kurva yang dua kali terdiferensialkan secara kontinu
\maabb {\gamma} {I} {M
} {}
pada suatu
\definitionsverweis {manifold Riemann}{}{}
$M$.
}

}
{

\aufzaehlungsechs{Misalkan $N$ suatu
\definitionsverweis {medan normal satuan}{}{}
pada $Y$ yang menentukan orientasi. Maka pemetaan
\maabbeledisp {L_P} {T_P Y} { T_PY
} {v} {- \left(  D_{v} N  \right) \left(  P  \right)
} {,}
disebut pemetaan Weingarten
pada $T_PY$.
}{Suatu
bundel vektor riil
$V$ ber-rank $r$ adalah ruang topologis yang dilengkapi dengan suatu
\definitionsverweis {pemetaan kontinu}{}{}
\maabb {p} {V} {X
} {}
sedemikian rupa sehingga setiap
\definitionsverweis {serat}{}{}
$p^{-1}(x)$ merupakan
\definitionsverweis {ruang vektor riil}{}{}
\definitionsverweis {berdimensi}{}{}
$r$, dan terdapat suatu
\definitionsverweis {penutup terbuka}{}{}

\mavergleichskette
{\vergleichskette
{X
}
{ = }{ \bigcup_{i \in I} U_i
}
{  }{
}
{    }{
}
{   }{
}
}
{}{}{}
serta
\definitionsverweis {homeomorfisme}{}{}
\maabbdisp {\varphi_i} {p^{-1}  \left(  U_i  \right) } { U_i \times \R^r
} {}
di atas $U_i$ yang pada setiap serat menginduksi suatu
\definitionsverweis {isomorfisme linear}{}{}
\maabbdisp {\left(  \varphi_i  \right)_x} {p^{-1} (x) } { \R^r
} {}
.
}{Integral lintasan didefinisikan oleh

\mavergleichskettedisp
{\vergleichskette
{  \int_\gamma \omega
}
{ =} {\int_a^b \gamma^* \omega
}
{ } {
}
{ } {
}
{ } {
}
}
{}{}{}
.
}{Suatu
$n$-\definitionsverweis {bentuk diferensial}{}{}
$\omega$ pada $M$ disebut bentuk volume positif jika, untuk setiap
\definitionsverweis {peta}{}{}
\maabbdisp {\alpha} { U } { V
} {}
\zusatzklammer {dengan \mathlk{V \subseteq \R^n}{} dan fungsi-fungsi koordinat \mathlk{x_1 , \ldots , x_n}{}} {} {}
dalam representasi lokal bentuk diferensial

\mavergleichskettedisp
{\vergleichskette
{ \alpha_* \omega
}
{ =} { f dx_1 \wedge \ldots \wedge dx_n
}
{ } {
}
{ } {
}
{ } {
}
}
{}{}{}
fungsi $f$ positif di setiap titik.
}{Suatu
\definitionsverweis {bentuk diferensial}{}{}
\definitionsverweis {diferensiabel}{}{}
$\omega$ pada $M$ disebut tertutup jika
\definitionsverweis {turunan luar}{}{}

\mathl{d \omega=0}{}.
}{Komposisi

\mathl{\pi_{\text{vert} }   \circ TT(\gamma)}{} disebut
percepatan tangensial
dari $\gamma$, dengan proyeksi tersebut berasal dari
\definitionsverweis {koneksi Levi-Civita}{}{.}
}

 }

__NOEDITSECTION__

\inputaufgabeklausurloesung
{3}

{

Rumuskan teorema-teorema berikut.
\aufzaehlungdrei{\stichwort {Teorema tentang kelengkungan normal} {.}}{Rumus untuk menghitung volume kanonik pada suatu manifold Riemann di dalam suatu peta.}{\stichwort {Teorema Stokes} {} untuk manifold dengan batas.}

}
{

\aufzaehlungdrei{\faktsituation {Misalkan

\mavergleichskette
{\vergleichskette
{ W
}
{ \subseteq }{ \R^n
}
{  }{
}
{    }{
}
{   }{
}
}
{}{}{}
\definitionsverweis {terbuka}{}{,}
\maabb {h} { W } { \R
} {}
suatu
\definitionsverweis {fungsi yang dua kali terdiferensialkan secara kontinu}{}{}
dan

\mavergleichskette
{\vergleichskette
{ Y
}
{ = }{ h^{-1}(c)
}
{  }{
}
{    }{
}
{   }{
}
}
{}{}{}
merupakan
\definitionsverweis {serat}{}{}
dari

\mavergleichskette
{\vergleichskette
{ c
}
{ \in }{ \R
}
{  }{
}
{    }{
}
{   }{
}
}
{}{}{,}
dengan $h$
\definitionsverweis {reguler}{}{}
di setiap titik $Y$. Misalkan

\mavergleichskette
{\vergleichskette
{ P
}
{ \in }{ Y
}
{  }{
}
{    }{
}
{   }{
}
}
{}{}{.}
Misalkan

\mavergleichskette
{\vergleichskette
{ E
}
{ \subseteq }{ \R^n
}
{  }{
}
{    }{
}
{   }{
}
}
{}{}{}
suatu
\definitionsverweis {bidang normal}{}{}
yang melalui $P$ pada $Y$.}
\faktfolgerung {Maka
\definitionsverweis {kelengkungan normal}{}{}
dari $Y$ di $P$ sama dengan
\definitionsverweis {kelengkungan}{}{}
kurva bidang

\mavergleichskette
{\vergleichskette
{ Y \cap (P+E)
}
{ \subseteq }{ P+E
}
{ \cong }{ \R^2
}
{    }{
}
{   }{
}
}
{}{}{}
di titik $P$.}
\faktzusatz {}
\faktzusatz {}}{Misalkan $M$ suatu
\definitionsverweis {manifold Riemann}{}{}
\definitionsverweis {berorientasi}{}{}
dan $\omega$
\definitionsverweis {bentuk volume kanonik}{}{.}
Misalkan
\maabbdisp {\alpha} { U } { V
} {}
suatu
\definitionsverweis {peta berorientasi}{}{}
dengan

\mavergleichskettedisp
{\vergleichskette
{ V
}
{ \subseteq} { \R^n
}
{ } {
}
{ } {
}
{ } {
}
}
{}{}{}
terbuka, dengan koordinat
\mathl{x_1 , \ldots , x_n}{} dan
\definitionsverweis {matriks fundamental metrik}{}{}

\mathl{G=(g_{ij})_{1 \leq i,j \leq n}}{}, serta
\mathl{g= \det  G}{.} Maka

\mavergleichskettedisp
{\vergleichskette
{  \alpha_* \omega
}
{ =} { \sqrt{ g} dx_1 \wedge \ldots \wedge dx_n
}
{ } {
}
{ } {
}
{ } {
}
}
{}{}{.}
Untuk suatu
\definitionsverweis {himpunan bagian terukur}{}{}

\mavergleichskette
{\vergleichskette
{ T
}
{ \subseteq }{ U
}
{  }{
}
{    }{
}
{   }{
}
}
{}{}{}
berlaku

\mavergleichskettedisp
{\vergleichskette
{
\int_{  T   }  \omega
}
{ =} {
\int_{ \alpha(T)  }   \sqrt{g} dx_1 \wedge \ldots \wedge dx_n
}
{ =} { \int_{ \alpha(T)  }   \sqrt{g}    \,  d \lambda^n
}
{ } {
}
{ } {
}
}
{}{}{.}}{Misalkan $M$ suatu
\definitionsverweis {manifold diferensiabel}{}{}
berdimensi $n$ yang
\definitionsverweis {berorientasi}{}{}
dan
\definitionsverweis {mempunyai batas}{}{}
$\partial M$, serta mempunyai
\definitionsverweis {basis topologi yang terhitung}{}{.}
Misalkan $\omega$ suatu
$(n-1)$-\definitionsverweis {bentuk diferensial}{}{}
yang
\definitionsverweis {terdiferensialkan secara kontinu}{}{}
dengan
\definitionsverweis {dukungan}{}{}
\definitionsverweis {kompak}{}{} pada $M$. Maka

\mavergleichskettedisp
{\vergleichskette
{
\int_{ M  }  d\omega
}
{ =} {
\int_{ \partial M  }  \omega
}
{ } {
}
{ } {
}
{ } {
}
}
{}{}{.}}

 }

__NOEDITSECTION__

\inputaufgabeklausurloesung
{0}

{

}
{  }

__NOEDITSECTION__

\inputaufgabeklausurloesung
{0}

{

}
{  }

__NOEDITSECTION__

\inputaufgabeklausurloesung
{4}

{

Buktikan teorema tentang hubungan antara kelengkungan dan vektor normal satuan dari suatu kurva bidang yang diparameterkan oleh panjang busur.

}
{

Pernyataan pertama langsung mengikuti definisi

\mavergleichskettedisp
{\vergleichskette
{ \kappa(t)
}
{ =} { \gamma_1'(t) \gamma_2^{\prime \prime} (t) - \gamma_2'(t) \gamma_1^{\prime \prime} (t)
}
{ } {
}
{ } {
}
{ } {
}
}
{}{}{.}
Dari

\mavergleichskettedisp
{\vergleichskette
{ \left(  \gamma_1'(t)  \right)^2 +  \left(  \gamma_2' (t)  \right)^2
}
{ =} { 1
}
{ } {
}
{ } {
}
{ } {
}
}
{}{}{}
diperoleh

\mavergleichskettedisp
{\vergleichskette
{ \gamma_1'(t) \gamma_1^{\prime \prime}(t) + \gamma_2'(t) \gamma_2^{\prime \prime} (t)
}
{ =} { 0
}
{ } {
}
{ } {
}
{ } {
}
}
{}{}{,}
karena itu $\gamma^{\prime \prime}(t)$ tegak lurus terhadap $\gamma'(t)$ dan bergantung linear pada
\mathl{N(t)}{.} Dengan demikian,

\mavergleichskettedisp
{\vergleichskette
{ \gamma^{\prime \prime}(t)
}
{ =} { \left\langle  \gamma^{\prime \prime}(t) ,  N(t)  \right\rangle  N(t)
}
{ =} { \kappa(t) N(t)
}
{ } {
}
{ } {
}
}
{}{}{}
dan oleh karena itu

\mavergleichskettedisp
{\vergleichskette
{ \Vert { \gamma^{\prime \prime}(t) } \Vert
}
{ =} { \betrag {  \kappa(t) }
}
{ } {
}
{ } {
}
{ } {
}
}
{}{}{.}

 }

__NOEDITSECTION__

\inputaufgabeklausurloesung
{5 (1+2+2)}

{

Pada
\definitionsverweis {permukaan bola satuan}{}{}
$Y$ berdimensi dua, kita meninjau
\definitionsverweis {medan vektor tangensial}{}{}

\mavergleichskettedisp
{\vergleichskette
{ F \begin{pmatrix}  x  \\ y \\ z   \end{pmatrix}
}
{ =} { \begin{pmatrix}  y  \\ -x\\  0   \end{pmatrix}
}
{ } {
}
{ } {
}
{ } {
}
}
{}{}{.}
\aufzaehlungdrei{Tentukan
\definitionsverweis {matriks Jacobi}{}{}
dari $F$ pada $\R^3$.
}{Untuk titik

\mavergleichskette
{\vergleichskette
{ P
}
{ = }{ \begin{pmatrix}  0  \\ 1 \\  0   \end{pmatrix}
}
{ \in }{ Y
}
{    }{
}
{   }{
}
}
{}{}{}
tentukan suatu matriks untuk pemetaan
\maabbeledisp {} {T_PY} { T_PY
} { v } { \nabla_v F
} {,}
terhadap suatu
\definitionsverweis {basis}{}{}
yang sesuai.
}{Untuk titik

\mavergleichskette
{\vergleichskette
{ Q
}
{ = }{ \begin{pmatrix}  0  \\ 0\\  1   \end{pmatrix}
}
{ \in }{ Y
}
{    }{
}
{   }{
}
}
{}{}{}
tentukan suatu matriks untuk pemetaan
\maabbeledisp {} {T_QY} { T_QY
} { v } { \nabla_v F
} {,}
terhadap suatu
\definitionsverweis {basis}{}{}
yang sesuai.
}

}
{

\aufzaehlungdrei{Matriks Jacobi-nya adalah

\mathdisp {\begin{pmatrix}  0   &   1  &  0  \\  -1 &  0  &  0 \\ 0  &  0 &  0  \end{pmatrix}} { . }

}{Sebagai medan normal, kita ambil $\begin{pmatrix}  x  \\ y \\ z   \end{pmatrix}$. Di

\mavergleichskette
{\vergleichskette
{ P
}
{ = }{ \begin{pmatrix}  0  \\ 1 \\  0   \end{pmatrix}
}
{  }{
}
{    }{
}
{   }{
}
}
{}{}{}
vektor-vektor
\mathkor {} {\begin{pmatrix}  1  \\ 0 \\  0   \end{pmatrix}} {dan} {\begin{pmatrix}  0  \\ 0\\  1   \end{pmatrix}} {}
membentuk suatu basis ruang singgung. Berlaku

\mavergleichskettedisp
{\vergleichskette
{ \begin{pmatrix}  0   &   1  &  0  \\  -1 &  0  &  0 \\ 0  &  0 &  0  \end{pmatrix} \begin{pmatrix}  1  \\ 0 \\  0   \end{pmatrix}
}
{ =} { \begin{pmatrix}  0  \\ -1\\  0   \end{pmatrix}
}
{ } {
}
{ } {
}
{ } {
}
}
{}{}{}
dan

\mavergleichskettedisp
{\vergleichskette
{  \begin{pmatrix}  0   &   1  &  0  \\  -1 &  0  &  0 \\ 0  &  0 &  0  \end{pmatrix} \begin{pmatrix}  0  \\ 0\\  1   \end{pmatrix}
}
{ =} { \begin{pmatrix}  0  \\ 0\\  0   \end{pmatrix}
}
{ } {
}
{ } {
}
{ } {
}
}
{}{}{.}
Di sini,
\mathl{\begin{pmatrix}  0  \\ -1\\  0   \end{pmatrix}}{} merupakan kelipatan vektor normal, sehingga proyeksi ortogonalnya pada ruang singgung juga merupakan vektor nol. Jadi, pemetaan
\mathl{v \mapsto \nabla_vF}{} adalah pemetaan nol.
}{Di

\mavergleichskette
{\vergleichskette
{ Q
}
{ = }{ \begin{pmatrix}  0  \\ 0\\  1   \end{pmatrix}
}
{  }{
}
{    }{
}
{   }{
}
}
{}{}{}
vektor-vektor
\mathkor {} {\begin{pmatrix}  1  \\ 0 \\  0   \end{pmatrix}} {dan} {\begin{pmatrix}  0  \\ 1 \\  0   \end{pmatrix}} {}
membentuk suatu basis ruang singgung. Berlaku

\mavergleichskettedisp
{\vergleichskette
{ \begin{pmatrix}  0   &   1  &  0  \\  -1 &  0  &  0 \\ 0  &  0 &  0  \end{pmatrix} \begin{pmatrix}  1  \\ 0 \\  0   \end{pmatrix}
}
{ =} { \begin{pmatrix}  0  \\ -1\\  0   \end{pmatrix}
}
{ } {
}
{ } {
}
{ } {
}
}
{}{}{}
dan

\mavergleichskettedisp
{\vergleichskette
{ \begin{pmatrix}  0   &   1  &  0  \\  -1 &  0  &  0 \\ 0  &  0 &  0  \end{pmatrix} \begin{pmatrix}  0  \\ 1 \\  0   \end{pmatrix}
}
{ =} { \begin{pmatrix}  1  \\ 0 \\  0   \end{pmatrix}
}
{ } {
}
{ } {
}
{ } {
}
}
{}{}{.}
Kedua vektor hasil tersebut sudah menyinggung $Y$ di $Q$. Suatu matriks yang merepresentasikan
\mathl{v \mapsto \nabla_vF}{} adalah

\mathdisp {\begin{pmatrix}  0   &   -1  \\  1  &  0 \end{pmatrix}} { . }

}

 }

__NOEDITSECTION__

\inputaufgabeklausurloesung
{4 (2+2)}

{

Kita meninjau elips

\mavergleichskettedisp
{\vergleichskette
{ E
}
{ =} { \left\{ (u,v)  \mid   2u^2 +5v^2 = 4         \right\}
}
{ } {
}
{ } {
}
{ } {
}
}
{}{}{.}
\aufzaehlungdrei{Definisikan suatu pemetaan surjektif yang terdiferensialkan secara kontinu
\maabb {} {\R} {E
} {.}
}{Deskripsikan suatu
\definitionsverweis {difeomorfisme}{}{}
antara lingkaran $S^1$ dan $E$.
}{
}

}
{

Kita menggunakan deskripsi

\mavergleichskettedisp
{\vergleichskette
{ E
}
{ =} { \left\{ (u,v)  \mid   { \frac{   1 }{  2 } } u^2+ { \frac{   5 }{  4 } } v^2 = 1         \right\}
}
{ } {
}
{ } {
}
{ } {
}
}
{}{}{.}
\aufzaehlungzwei {Pemetaan
\maabbeledisp {} {\R} { E
} { \alpha} { \left(  \sqrt{2} \cos  \alpha   , \,   { \frac{   2 }{  \sqrt{5}  } } \sin  \alpha                   \right)
} {}
terdiferensialkan secara kontinu sebagai pemetaan ke $\R^2$, dan citranya terletak pada $E$. Oleh karena itu, pemetaan ini juga terdiferensialkan secara kontinu sebagai pemetaan ke $E$. Surjektivitasnya sekaligus dibuktikan dalam (2).
} {Kita meninjau pemetaan
\maabbeledisp {} {S^1} {E
} {(x,y)} { \left(  \sqrt{2} x  , \,   { \frac{   2 }{  \sqrt{5}  } } y                   \right)
} {.}
Sebagai pembatasan suatu pemetaan linear bijektif pada $\R^2$, pemetaan ini terdiferensialkan secara kontinu dan injektif, serta memang bernilai di $E$
\zusatzklammer {berdasarkan persamaan-persamaan eksplisit} {} {.}
Pemetaan ini juga surjektif karena kita dapat langsung memberikan pemetaan invers
\maabbeledisp {} { E } {S^1
} {(u,v)} { \left(  { \frac{   1 }{  \sqrt{2}  } } u  , \,   { \frac{   \sqrt{5}  }{  2  } } v                   \right)
} {,}
secara eksplisit. Karena parametrisasi trigonometrik lingkaran satuan bersifat surjektif, pemetaan yang diberikan dalam (1) juga surjektif.
}

 }

__NOEDITSECTION__

\inputaufgabeklausurloesung
{0}

{

}
{  }

__NOEDITSECTION__

\inputaufgabeklausurloesung
{4}

{

Berikan suatu
\definitionsverweis {medan vektor}{}{} yang
\definitionsverweis {kontinu}{}{}
pada $S^2$ dan hanya mempunyai satu titik nol.

}
{

Kita meninjau pada $\R^2$ medan vektor kontinu $F$ yang diberikan oleh

\mavergleichskettedisp
{\vergleichskette
{F(x,y)
}
{ =} { { \frac{   1 }{  x^2+y^2 } }  e_1
}
{ } {
}
{ } {
}
{ } {
}
}
{}{}{}
Medan ini tidak mempunyai titik nol dan bersifat kontinu. Kita mengangkut medan vektor ini melalui proyeksi stereografis ke

\mavergleichskette
{\vergleichskette
{  \R^2
}
{ \cong }{ S^2 \setminus \{N\}
}
{  }{
}
{    }{
}
{   }{
}
}
{}{}{}
dan memperluasnya di Kutub Utara dengan nilai

\mavergleichskette
{\vergleichskette
{ 0
}
{ \in }{  T_NS^2
}
{  }{
}
{    }{
}
{   }{
}
}
{}{}{.}
Kita menyatakan bahwa medan vektor ini kontinu. Untuk itu, misalkan
\mathl{P_n}{} suatu barisan pada $S^2$ yang konvergen ke $N$. Kita dapat langsung menganggap bahwa semuanya memenuhi

\mavergleichskette
{\vergleichskette
{P_n
}
{ \neq }{N
}
{  }{
}
{    }{
}
{   }{
}
}
{}{}{}
. Citra peta barisan ini adalah

\mavergleichskettedisp
{\vergleichskette
{P'_n
}
{ =} { (x_n,y_n)
}
{ } {
}
{ } {
}
{ } {
}
}
{}{}{,}
dan karena $P_n$ konvergen ke Kutub Utara,
\mathl{\betrag { P_n' }}{} divergen ke $\infty$. Oleh karena itu, barisan

\mavergleichskettedisp
{\vergleichskette
{ F(P'_n)
}
{ =} { F(x_n , y_n)
}
{ =} { { \frac{   1 }{  x^2+y^2 } } e_1
}
{ } {
}
{ } {
}
}
{}{}{}
konvergen ke $0$.

 }

__NOEDITSECTION__

\inputaufgabeklausurloesung
{6}

{

Misalkan diberikan suatu
\definitionsverweis {data perekatan}{}{}

\mathbed {U_i} {}
 {i \in I} {}
 {} {} {} {,}
untuk
\definitionsverweis {ruang-ruang topologis}{}{.}
Buktikan bahwa terdapat suatu ruang topologis $X$ yang ditentukan secara unik, suatu penutup terbuka

\mavergleichskette
{\vergleichskette
{ X
}
{ = }{ \bigcup_{i \in I} V_i
}
{  }{
}
{    }{
}
{   }{
}
}
{}{}{}
dan
\definitionsverweis {homeomorfisme}{}{}
\maabb {\psi_i} {U_i } { V_i
} {}
sedemikian sehingga

\mavergleichskettedisp
{\vergleichskette
{ \psi_i  \left(  U_{ij}  \right)
}
{ =} { V_i \cap V_j
}
{ } {
}
{ } {
}
{ } {
}
}
{}{}{}
dan

\mavergleichskettedisp
{\vergleichskette
{ \psi_i \vert_{U_{ij} }
}
{ =} { \psi_j \vert_{U_{ji} }  \circ   \varphi_{ji}
}
{ } {
}
{ } {
}
{ } {
}
}
{}{}{}.

}
{

Misalkan $Y$ gabungan saling lepas dari himpunan-himpunan $U_i$. Pada $Y$, kita mendefinisikan suatu
\definitionsverweis {relasi ekuivalensi}{}{}
$\sim$ dengan menganggap titik-titik

\mavergleichskette
{\vergleichskette
{ x_i
}
{ \in }{ U_i
}
{  }{
}
{    }{
}
{   }{
}
}
{}{}{}
dan

\mavergleichskette
{\vergleichskette
{ x_j
}
{ \in }{ U_j
}
{  }{
}
{    }{
}
{   }{
}
}
{}{}{}
ekuivalen jika

\mavergleichskette
{\vergleichskette
{ x_i
}
{ \in }{ U_{ij}
}
{  }{
}
{    }{
}
{   }{
}
}
{}{}{,}

\mavergleichskette
{\vergleichskette
{ x_j
}
{ \in }{ U_{ji}
}
{  }{
}
{    }{
}
{   }{
}
}
{}{}{}
dan

\mavergleichskette
{\vergleichskette
{ \varphi_{ji} (x_i)
}
{ = }{ x_j
}
{  }{
}
{    }{
}
{   }{
}
}
{}{}{}
berlaku. Sifat-sifat relasi ekuivalensi dijamin oleh syarat kosikel; lihat
[[Topologischer Raum/Verklebungsdatum/Äquivalenzrelation/Aufgabe|[[:Kurs:Bündel, Garben und Kohomologie (Osnabrück 2019-2020)/4/Klausur mit Lösungen/latex]] (Bündel, Garben und Kohomologie (Osnabrück 2019-2020))[[Kategorie:Kurs:Differentialgeometrie/Referenznummer|Bündel, Garben und Kohomologie (Osnabrück 2019-2020)]] (Differentialgeometrie (Osnabrück 2023))]].
Kita menetapkan

\mavergleichskettedisp
{\vergleichskette
{X
}
{ \defeq} { Y/ \sim
}
{ } {
}
{ } {
}
{ } {
}
}
{}{}{}
dan memperlengkapi $X$ dengan
\definitionsverweis {topologi hasil bagi}{}{.}
Komposisi-komposisi

\mathdisp {U_i \longrightarrow Y \longrightarrow X} {  }

adalah $\psi_i$, sedangkan $V_i$ merupakan citra pemetaan-pemetaan tersebut. Dengan demikian, diperoleh homeomorfisme
\maabb {\psi_i} {U_i } { V_i
} {}
. Untuk

\mavergleichskette
{\vergleichskette
{ x
}
{ \in }{ U_i
}
{  }{
}
{    }{
}
{   }{
}
}
{}{}{}

\mavergleichskettedisp
{\vergleichskette
{ \psi_i (x)
}
{ \in} { V_j
}
{ } {
}
{ } {
}
{ } {
}
}
{}{}{}
berlaku tepat jika

\mavergleichskette
{\vergleichskette
{ x
}
{ \in }{ U_{ij}
}
{  }{
}
{    }{
}
{   }{
}
}
{}{}{}
karena tepat dalam hal ini $x$ diidentifikasi dengan
\mathl{\varphi_{ij}(x)}{}. Oleh karena itu,

\mavergleichskette
{\vergleichskette
{ \psi_i(U_{ij})
}
{ = }{ V_i \cap V_j
}
{  }{
}
{    }{
}
{   }{
}
}
{}{}{.}
Komutativitas diagram

\mathdisp {\begin{matrix} U_{ij}   & \stackrel{ \varphi_{ji} }{\longrightarrow} &  U_{ji}   & \\ & \!\!\! \!\! \psi_i \searrow & \downarrow \psi_j \!\!\! \!\! & \\ & &  V_i \cap V_j   & \!\!\!\!\! \!\!\!  \\ \end{matrix}} {  }

juga diperoleh.

 }

__NOEDITSECTION__

\inputaufgabeklausurloesung
{2}

{

Misalkan $M$ suatu
\definitionsverweis {manifold diferensiabel}{}{,}
\maabb {\gamma} { I } { M
} {}
suatu
\definitionsverweis {kurva diferensiabel}{}{}
di $M$, dan $\omega$ suatu
$1$-\definitionsverweis {bentuk}{}{}
terdiferensialkan pada $M$. Buktikan bahwa

\mavergleichskettedisp
{\vergleichskette
{ \gamma^* (d \omega)
}
{ =} { 0
}
{ } {
}
{ } {
}
{ } {
}
}
{}{}{}.

}
{

$d \omega$ merupakan suatu bentuk-$2$ pada $M$, dan sifat ini terbawa ke $\gamma^*(d \omega)$. Akan tetapi, pada manifold berdimensi satu, semua bentuk-$2$ sama dengan $0$.

 }

__NOEDITSECTION__

\inputaufgabeklausurloesung
{6 (3+3)}

{

Misalkan $M$ suatu
\definitionsverweis {manifold diferensiabel}{}{}
yang
\definitionsverweis {kompak}{}{}
dan mempunyai sedikitnya dua titik.
\aufzaehlungzwei {Buktikan bahwa suatu
$1$-\definitionsverweis {bentuk diferensial}{}{}
kontinu yang
\definitionsverweis {eksak}{}{}
pada $M$ mempunyai sedikitnya dua titik nol.
} {Buktikan bahwa hal ini tidak harus berlaku bagi suatu bentuk diferensial yang
\definitionsverweis {tertutup}{}{.}
}

}
{

\aufzaehlungzwei {Suatu bentuk diferensial eksak dapat dituliskan sebagai

\mavergleichskettedisp
{\vergleichskette
{ \omega
}
{ =} { df
}
{ } {
}
{ } {
}
{ } {
}
}
{}{}{}
dengan suatu fungsi diferensiabel
\maabbdisp {f} { M } {\R
} {}
. Jika $f$ konstan, maka

\mavergleichskette
{\vergleichskette
{ df
}
{ = }{ 0
}
{  }{
}
{    }{
}
{   }{
}
}
{}{}{}
dan pernyataannya jelas. Jadi, misalkan $f$ tidak konstan. Menurut
[[Kompaktheit/Stetige Funktion/Maximum wird angenommen/Fakt|Fakt]] [[Kurs:Differentialgeometrie/Kompaktheit/Stetige Funktion/Maximum wird angenommen/Fakt/Faktreferenznummer|*****]]
fungsi kontinu tersebut mencapai maksimum global dan minimum global, katakanlah masing-masing di titik
\mathkor {} {P} {dan} {Q} {.}
Menurut
[[Differenzierbare Mannigfaltigkeit/Funktion/Extrema/Differentialform/Nullstelle/Aufgabe|Aufgabe]] [[Kurs:Differentialgeometrie/Differenzierbare Mannigfaltigkeit/Funktion/Extrema/Differentialform/Nullstelle/Aufgabe/Aufgabereferenznummer|*****]]
maka

\mavergleichskettedisp
{\vergleichskette
{ (df)_P
}
{ =} { (df)_Q
}
{ =} { 0
}
{ } {
}
{ } {
}
}
{}{}{.}
} {Pada lingkaran berdimensi satu yang kompak $S^1$, kita meninjau bentuk diferensial

\mavergleichskettedisp
{\vergleichskette
{ \omega(x,y)
}
{ =} {  ydx - xdy
}
{ } {
}
{ } {
}
{ } {
}
}
{}{}{.}
Ini memang merupakan suatu bentuk diferensial pada lingkaran karena
\mathl{\left(  y   , \,   -x                   \right)}{} bersifat tangensial terhadap vektor posisi
\mathl{(x,y)}{}. Bentuk ini tertutup karena dimensinya $1$. Bentuk ini tidak mempunyai titik nol karena pada lingkaran,
\mathkor {} {x} {dan} {y} {}
tidak mungkin sekaligus bernilai $0$.
}

 }

__NOEDITSECTION__

\inputaufgabeklausurloesung
{4}

{

Berikan suatu contoh
\definitionsverweis {manifold dengan batas}{}{,}
yang
\definitionsverweis {terhubung}{}{}
dan batasnya terdiri atas
\zusatzklammer {gabungan saling lepas dari} {} {}
tak terhingga banyak salinan $\R^1$.

}
{

Fungsi
\maabbeledisp {h} {]-1,1[ } { \R
} { x } { { \frac{   1  }{  x^2-1 } }
} {,}
menuju tak hingga pada kedua ujungnya. Grafik $h$ difeomorf dengan $]-1,1[$, dan dengan demikian juga dengan $\R$. Sekarang, untuk

\mavergleichskette
{\vergleichskette
{ n
}
{ \in }{  \Z
}
{  }{
}
{    }{
}
{   }{
}
}
{}{}{}
, kita meninjau fungsi yang bersesuaian dengan domain yang digeser sebesar $4n$
\zusatzklammer {jadi dengan domain
\mathl{]4n-1,4n+1[}{}} {} {.}
Misalkan $G_n$ grafik yang bersesuaian dan $U_n$ epigraf terbuka
\zusatzklammer {!} {} {}
yang bersesuaian. Kita menetapkan

\mavergleichskettedisp
{\vergleichskette
{ M
}
{ =} { \R^2 \setminus \biguplus_{n \in \Z} U_n
}
{ } {
}
{ } {
}
{ } {
}
}
{}{}{.}
Himpunan ini merupakan suatu manifold terhubung dengan batas, dan batasnya adalah gabungan saling lepas dari himpunan-himpunan $G_n$.

 }

__NOEDITSECTION__

\inputaufgabeklausurloesung
{10}

{

Buktikan \stichwort {teorema tentang partisi kesatuan} {.}

}
{

Menurut
[[Mannigfaltigkeit/Abzählbare Basis/Lokal endlicher Atlas/Fakt|Lemma 22.9 (Differentialgeometrie (Osnabrück 2023))]]
kita dapat menganggap bahwa terdapat suatu penutup terbuka oleh daerah-daerah peta

\mathbed {U_j} {}
 {j \in J} {}
 {} {} {} {,}
\zusatzklammer {$J$ terhitung} {} {}
dengan
\maabbdisp {\alpha_j} {U_j } { V_j
} {}
dan lingkungan-lingkungan bola

\mavergleichskettedisp
{\vergleichskette
{ U \left(   0 , \delta_j   \right)
}
{ \subset} { B  \left(  0 , \epsilon_j  \right)
}
{ \subset} { V_j
}
{ } {
}
{ } {
}
}
{}{}{}
\zusatzklammer {dengan

\mavergleichskettek
{\vergleichskettek
{ \delta_j
}
{ < }{ \epsilon_j
}
{  }{
}
{    }{
}
{   }{
}
}
{}{}{}} {} {}
sedemikian rupa sehingga himpunan-himpunan
\mathl{\alpha_j^{-1} \left(  U \left(   0 , \delta_j   \right)  \right)}{} juga membentuk suatu penutup dari $M$, dan setiap titik

\mavergleichskette
{\vergleichskette
{ P
}
{ \in }{ M
}
{  }{
}
{    }{
}
{   }{
}
}
{}{}{}
hanya termuat dalam berhingga banyak himpunan $U_j$, dan khususnya hanya dalam berhingga banyak himpunan
\mathl{\alpha_j^{-1} \left(  U \left(   0 , \delta_j   \right)  \right)}{}. Pada $V_j$, kita meninjau fungsi $g_j$ yang diberikan oleh

\mavergleichskettedisp
{\vergleichskette
{ g_j (v)
}
{ =} { \begin{cases}  e^{ - { \frac{   1  }{  \left(  \delta_j^2- \Vert { v } \Vert^2  \right)^2 } } }  \text{ jika } \Vert { v } \Vert < \delta_j \, , \\ 0 \text{ selain itu} \, , \end{cases}
}
{ } {
}
{ } {
}
{ } {
}
}
{}{}{}
. Fungsi ini bernilai positif tepat pada
\mathl{U \left(   0 , \delta_j   \right)}{}, dan
\definitionsverweis {dukungannya}{}{}
adalah
\mathl{B  \left(  0 , \delta_j  \right)}{.} Dengan meninjau dua himpunan bagian terbuka
\zusatzklammer {yang menutupi $V_j$} {} {}
\mathkor {} {U \left(   0 , \epsilon_j   \right)} {dan} {V_j \setminus B  \left(  0 , \delta_j  \right)} {}
diperoleh bahwa $g_j$ diferensiabel tak hingga kali. Kita mendefinisikan suatu fungsi
\maabbdisp {\tilde{g}_j} { M } {\R
} {}
melalui

\mavergleichskettedisp
{\vergleichskette
{ \tilde{g}_j(x)
}
{ =} { \begin{cases}  g( \alpha_j(x)) \text{ jika }  x \in \alpha_j^{-1}( U \left(   0 , \delta_j   \right) ) \, , \\  0 \text{ selain itu} \, . \end{cases}
}
{ } {
}
{ } {
}
{ } {
}
}
{}{}{}
Fungsi ini
\definitionsverweis {terdiferensialkan secara kontinu}{}{}
pada $M$ karena \anfuehrung{pita}{}
\mathl{B  \left(  0 , \epsilon_j  \right) \setminus U \left(   0 , \delta_j   \right)}{} memungkinkan transisi mulus. Kita menetapkan

\mavergleichskettedisp
{\vergleichskette
{ \tilde{g}(x)
}
{ \defeq} { \sum_{j \in J} \tilde{g}_j(x)
}
{ } {
}
{ } {
}
{ } {
}
}
{}{}{,}
yang merupakan jumlah berhingga pada setiap titik karena
\definitionsverweis {dukungan}{}{}
dari $\tilde{g}_j$ terletak di dalam

\mavergleichskettedisp
{\vergleichskette
{ \alpha^{-1} \left(  B  \left(  0 , \delta_j  \right)  \right)
}
{ \subset} { U_j
}
{ } {
}
{ } {
}
{ } {
}
}
{}{}{}
. Fungsi ini
\definitionsverweis {terdiferensialkan secara kontinu}{}{}
pada $M$ dan positif di setiap titik karena
\mathl{\tilde{g}_j(x)}{} positif pada himpunan-himpunan penutup
\mathl{\alpha^{-1} \left(  U \left(   0 , \delta_j   \right)  \right)}{}. Maka fungsi-fungsi

\mavergleichskettedisp
{\vergleichskette
{ h_j
}
{ =} { { \frac{   \tilde{g}_j  }{ \tilde{g}  } }
}
{ } {
}
{ } {
}
{ } {
}
}
{}{}{}
membentuk partisi kesatuan yang dicari.

 }

__NOEDITSECTION__

\inputaufgabeklausurloesung
{10 (2+6+2)}

{

Kita meninjau
$2$-\definitionsverweis {bentuk diferensial}{}{}

\mavergleichskettedisp
{\vergleichskette
{ \omega
}
{ =} { x dy \wedge dz -y dx \wedge dz + z dx \wedge dy
}
{ } {
}
{ } {
}
{ } {
}
}
{}{}{}
pada bola satuan
\mathl{B  \left(  0 , 1  \right)}{.}
\aufzaehlungdreiabc{Buktikan bahwa $d \omega$ adalah tiga kali bentuk volume standar pada bola tersebut.
}{Buktikan bahwa
\mathl{\omega\vert_{S^2}}{} merupakan bentuk luas permukaan standar pada permukaan bola satuan.
}{Hitung luas permukaan bola dari
[[Allgemeines Kugelvolumen/Mit Cavalieri-Prinzip/Beispiel|volume bola]]
dengan menggunakan
[[Mannigfaltigkeit mit Rand/Satz von Stokes/Fakt|Teorema Stokes]].
}

}
{

a) Berlaku

\mavergleichskettealign
{\vergleichskettealign
{d \omega
}
{ =} { d \left(  x dy \wedge dz  -y dx \wedge dz +  z dx \wedge dy \right)
}
{ =} { dx \wedge dy \wedge dz - dy \wedge dx \wedge dz + d z \wedge dx \wedge dy
}
{ =} { dx \wedge dy \wedge dz + dx \wedge dy \wedge dz + d x \wedge dy \wedge dz
}
{ =} { 3 d x \wedge dy \wedge dz
}
}
{}
{}{,}
di mana kita menggunakan fakta bahwa tanda berubah ketika faktor-faktor dalam produk baji dipertukarkan.

b) Untuk suatu titik
\mathl{P = \begin{pmatrix}  x  \\ y \\  z   \end{pmatrix}  \in S^2}{} dan dua vektor singgung ortonormal
\zusatzklammer {yang bersama dengan $P$ merepresentasikan orientasi} {} {}
\mathl{u= \begin{pmatrix} u_1  \\u_2 \\ u_3   \end{pmatrix}}{} serta
\mathl{v= \begin{pmatrix} v_1  \\v_2 \\ v_3   \end{pmatrix}}{}, menurut
[[Dachprodukt/Natürliche Dualität/Endlichdimensional/Fakt|Satz 58.4 (Lineare Algebra (Osnabrück 2024-2025))]]

\mavergleichskettealign
{\vergleichskettealign
{ \omega \vert_{S^2}  (u,v)
}
{ =} { \omega (u,v)
}
{ =} { z \det  \begin{pmatrix} u_1   &  v_1   \\ u_2  & v_2  \end{pmatrix} - y \det  \begin{pmatrix} u_1   &  v_1   \\ u_3  & v_3  \end{pmatrix} + x \det  \begin{pmatrix} u_2   &  v_2   \\ u_3  & v_3  \end{pmatrix}
}
{ =} { \det  \begin{pmatrix}  x   &  u_1  & v_1  \\  y  & u_2  & v_2  \\ z  & u_3  & v_3  \end{pmatrix}
}
{ } {
}
}
{}
{}{.}
Determinan suatu sistem ortonormal yang merepresentasikan orientasi adalah $1$. Jadi, bentuk diferensial yang dibatasi tersebut memenuhi sifat yang mengarakterisasi
\definitionsverweis {bentuk volume standar}{}{.}

c) Dengan menggunakan a), b), Teorema Stokes, dan volume bola, luas permukaan bola adalah

\mavergleichskettedisp
{\vergleichskette
{ \int_{S^2} \omega
}
{ =} { \int_{ B  \left(  0, 1 \right)  }  d \omega
}
{ =} { \int_{ B  \left(  0, 1 \right)  } 3 d x \wedge dy \wedge dz
}
{ =} { 3 \lambda^3 (   B  \left(  0, 1 \right)     )
}
{ =} { 4 \pi
}
}
{}{}{.}

 }

__NOEDITSECTION__

\inputaufgabeklausurloesung
{0}

{

}
{  }

 [[Kategorie:Latexseite]]
